% !TEX program = xelatex
% =========================================================
%  2025--2026《仿生算法的数学基础》课程考察 · 数据分析报告
%  基于 GA / ACO / PSO 三种仿生算法的对比分析
%  ——以教材 Appendix C 测试函数与 TSPLIB 真实 TSP 为例
% =========================================================
\documentclass[12pt,a4paper,fontset=windows]{ctexart}

% --- 几何与排版 ---
\usepackage[margin=2.5cm]{geometry}
\usepackage{setspace}
\onehalfspacing
\usepackage{indentfirst}
\setlength{\parskip}{4pt}

% --- 数学 ---
\usepackage{amsmath,amssymb,amsthm,mathtools,bm}

% --- 图与三线表 ---
\usepackage{graphicx}
\usepackage{float}
\usepackage{caption}
\captionsetup{font=small,labelfont=normalfont,labelsep=quad}
\usepackage{booktabs}
\usepackage{array}
\usepackage{tabularx}

% --- 超链接：全黑、无下划线、无色 ---
\usepackage[hidelinks]{hyperref}

% --- 流程图 ---
\usepackage{tikz}
\usetikzlibrary{shapes.geometric,arrows.meta,positioning,calc}

% --- 图片搜索路径 ---
\graphicspath{{../figures/}}

% --- 标题样式：中文用黑体字族（字体替换、非 \bfseries 加粗修饰）---
\ctexset{
  section/format={\Large\heiti\centering},
  subsection/format={\large\heiti},
  subsubsection/format={\normalsize\heiti},
  paragraph/format={\normalsize\heiti},
}
\setcounter{secnumdepth}{3}

% --- 数学算子与简写 ---
\newcommand{\R}{\mathbb{R}}
\DeclareMathOperator*{\argmin}{arg\,min}
\DeclareMathOperator*{\argmax}{arg\,max}

% --- 页面样式：去掉顶部章节页眉，仅保留底部居中页码 ---
\pagestyle{plain}

% =========================================================
\begin{document}

% ----------------- 封面 -----------------
\begin{titlepage}
  \centering
  \vspace*{2.5cm}
  {\Huge\heiti 基于 GA / ACO / PSO 三种仿生算法的\\[8pt]
   对比分析报告\par}
  \vspace{0.6cm}

  \vspace{2.4cm}

  \begin{tabular}{rl}
    \heiti 学\quad 号: & \underline{\hspace{4.5cm}} \\[8pt]
    \heiti 姓\quad 名: & \underline{\hspace{4.5cm}} \\[8pt]
  \end{tabular}
\end{titlepage}

\tableofcontents
\newpage

% =========================================================
\section{数据集来源及理解}
\label{sec:data}

本报告同时使用仿真数据与真实数据两类数据集，覆盖连续优化与组合优化两种典型场景，
用于全面评估GA、ACO、PSO三种仿生算法的精度、收敛速度、稳健性及其适用边界。

\subsection{仿真数据：教材 Appendix C测试函数}

仿真部分选用Evolutionary Optimization AlgorithmsAppendix C中的5个
经典无约束连续测试函数，覆盖单峰、谷型、规则强多峰、多峰带浅局部极小、多峰弱耦合等多种
地形难度。所有函数均为最小化问题且全局最优值 $f^{\ast}=0$。函数表达式、教材规定的搜索域
与全局最优点见表~\ref{tab:func}。

\begin{table}[H]
  \centering
  \caption{Appendix C 5个测试函数（维度 $d=30$，最小化）}
  \label{tab:func}
  \small
  \setlength{\tabcolsep}{4pt}
  \begin{tabular}{lllll}
    \toprule
    名称 & 教材节号 & 表达式 & 定义域（各维） & 全局最优点 \\
    \midrule
    Sphere     & C.1.1  & $\sum_i x_i^2$                                                                  & $[-5.12,\,5.12]$   & $\bm{x}=\bm{0}$ \\
    Rosenbrock & C.1.4  & $\sum_i 100(x_{i+1}\!-\!x_i^2)^2+(x_i\!-\!1)^2$                                & $[-2.048,\,2.048]$ & $\bm{x}=\bm{1}$ \\
    Rastrigin  & C.1.11 & $10d+\sum_i[x_i^2-10\cos(2\pi x_i)]$                                            & $[-5.12,\,5.12]$   & $\bm{x}=\bm{0}$ \\
    Ackley     & C.1.2  & $-20\,e^{-0.2\sqrt{\sum x_i^2/d}} - e^{\sum\cos(2\pi x_i)/d} + 20 + e$         & $[-30,\,30]$       & $\bm{x}=\bm{0}$ \\
    Griewank   & C.1.6  & $\sum_i x_i^2/4000 - \prod_i\cos(x_i/\sqrt{i}) + 1$                            & $[-600,\,600]$     & $\bm{x}=\bm{0}$ \\
    \bottomrule
  \end{tabular}
\end{table}

\subsection{真实数据：TSPLIB真实城市坐标TSP}

真实数据采用 TSPLIB（Reinelt 1991）—— 求解旅行商问题（TSP）的国际标准基准库，
坐标为真实地理或电路布局位置。本报告选取 4 个不同规模的实例，均采用 EUC\_2D 边权
$d_{ij}=\mathrm{round}\!\big(\!\sqrt{\Delta x^2+\Delta y^2}\big)$，并配以已知最优巡回长度
用于评估各算法的最优性差距 gap\%。实例见表~\ref{tab:tsp}。

\begin{table}[H]
  \centering
  \caption{TSPLIB 实例及已知最优巡回长度}
  \label{tab:tsp}
  \small
  \begin{tabular}{lrlr}
    \toprule
    实例 & 城市数 $n$ & 数据来源 & 已知最优 \\
    \midrule
    berlin52 & 52  & 柏林 52 个真实地点（Groetschel）   & 7542  \\
    eil51    & 51  & 51 城测试集（Christofides--Eilon）  & 426   \\
    kroA100  & 100 & 100 城（Krolak）                    & 21282 \\
    ch150    & 150 & 150 城（Churritz）                  & 6528  \\
    \bottomrule
  \end{tabular}
\end{table}

% =========================================================
\section{问题描述与解决思路}
\label{sec:formulation}

\subsection{分析目标}

\paragraph{总目标}
在连续优化与组合优化两类问题上，系统比较 GA、ACO、PSO 三种仿生算法的最终精度、
收敛速度与重复实验稳健性，回答两个核心问题：

\begin{enumerate}
  \item 不同算法在不同问题维度、不同地形特性、不同编码空间下的相对优劣是什么？
  \item 同一算法思想（如 ACO）在其本职问题域与外推到其他问题域时，性能差异是否显著？
\end{enumerate}

\paragraph{仿真数据：连续优化的数学化}
对每个测试函数 $f:\R^d\to\R$（见表~\ref{tab:func}）求
\begin{equation}
  \min_{\bm{x}\in[lb,ub]^d}\ f(\bm{x}),
\end{equation}
评价指标采用最终最优值 $f_{\rm best}$、达成阈值 $f<10^{-2}$ 的成功率，以及按函数评估次数
（FE）记录的 best-so-far 收敛曲线。

\paragraph{真实数据：TSP 的数学化}
给定 $n$ 个城市的坐标集合，求访问每个城市恰好一次并回到起点的最短闭合回路
$\pi=(\pi_1,\dots,\pi_n)\in S_n$：
\begin{equation}
  \min_{\pi}\ L(\pi) = \sum_{i=1}^{n-1} d(\pi_i,\pi_{i+1}) + d(\pi_n, \pi_1).
\end{equation}
评价指标采用最终巡回长度 $L_{\rm best}$、相对已知最优的差距
\[
  \mathrm{gap}\,\%=\frac{L_{\rm best}-L^{\ast}}{L^{\ast}}\times 100\%,
\]
以及收敛曲线与最优路线图。

\subsection{具体方法和流程}

\paragraph{算法形态对照}
GA、ACO、PSO 的核心思想（重组与变异、信息素引导、群体协同）可以同时实例化到连续空间与
排列空间，但所需算子完全不同。本报告对三种算法在两类问题上分别采用如下两种形态：

\begin{itemize}
  \item \emph{GA}：连续测试函数上采用实数编码，算子为[锦标赛选择 + SBX 交叉 + 多项式变异 + 精英保留]；
        TSP 上改为排列编码，算子为[锦标赛选择 + 顺序交叉 OX + 反转变异 + 精英保留]。
  \item \emph{ACO}：连续测试函数上采用 \emph{ACOR}（Socha \& Dorigo 2008）—— 维护大小为 $k$
        的解档案，逐维按 rank 权重选取引导解并以高斯核采样新解，模拟信息素的连续表征；
        TSP 上采用经典 \emph{Ant System}（Dorigo 1996）—— 信息素矩阵 $\tau$ 与启发式 $\eta=1/d_{ij}$
        协同引导蚂蚁概率构造巡回，迭代蒸发与沉积。
  \item \emph{PSO}：连续测试函数上采用标准 gbest PSO —— 线性衰减的惯性权重 $w$、速度钳制、
        个体最优与全局最优的双重引导；TSP 上改为离散交换序列 PSO（Clerc 2004）——
        把"速度"重新定义为有序交换算子序列，按惯性、认知、社会三个概率把"位置变成
        $p_{\rm best}$ / $g_{\rm best}$ 所需的交换序列"施加到当前排列上。
\end{itemize}

\paragraph{选型理由}
ACO 由 Dorigo 等人最初为求解 TSP 而提出，将其置于 TSP 上是其\emph{本职问题}，
符合"在合适问题上对比合适算法"的方法论；而要在连续测试函数上保持三算法齐全的公平对比，
ACO 需要采用其连续推广 ACOR（Socha \& Dorigo 2008），利用解档案与高斯核采样模拟
信息素的连续表征。GA 与 PSO 均同时支持实数与排列编码，分别采用各自标准的实数与离散形态。

\paragraph{公平性约定}
为保证横向可比，三种算法在同一问题、同一随机种子下满足如下约束：

\begin{itemize}
  \item 相同的函数评估（FE）预算：连续 30000、TSP 20000；
  \item 相互独立的随机数发生器，由共享 seed 派生，问题随机性与算法随机性互不耦合；
  \item 算法只通过问题对象的 \emph{evaluate} 接口访问目标，FE 计数、best-so-far 曲线与预算
        判定由问题对象内部统一维护，算法不自报指标。
\end{itemize}

\paragraph{实验流程}
图~\ref{fig:pipeline} 给出本报告的实验执行流程：所有算法在固定 FE 预算内循环
"生成–评估–更新"三步骤直至预算耗尽，最后返回 best-so-far 解与收敛曲线。

\begin{figure}[H]
  \centering
  \begin{tikzpicture}[
    node distance=10mm,
    box/.style={draw,rectangle,rounded corners=2pt,
                minimum width=38mm,minimum height=9mm,align=center,font=\small},
    decision/.style={draw,diamond,aspect=2,inner sep=0pt,
                     align=center,font=\small,minimum width=34mm,minimum height=12mm},
    arr/.style={-Latex,thick},
  ]
    \node[box]                              (init) {初始化种群 / 信息素 / 解档案};
    \node[box, below=of init]               (gen)  {生成新候选解};
    \node[box, below=of gen]                (eval) {评估目标值并累计 FE};
    \node[box, below=of eval]               (upd)  {更新历史最优 / 信息素 / 档案};
    \node[decision, below=of upd]           (cond) {FE 预算耗尽？};
    \node[box, below=of cond]               (ret)  {输出最优解与收敛曲线};
    \draw[arr] (init) -- (gen);
    \draw[arr] (gen)  -- (eval);
    \draw[arr] (eval) -- (upd);
    \draw[arr] (upd)  -- (cond);
    \draw[arr] (cond) -- node[right,font=\small,pos=0.4]{是} (ret);
    \coordinate (loop) at ([xshift=2.2cm]cond.east);
    \draw[arr] (cond.east) -- node[above,font=\small,pos=0.55]{否} (loop)
                |- (gen.east);
  \end{tikzpicture}
  \caption{实验执行流程图：所有算法共用同一主循环结构，"生成新候选解"步骤由各算法自身的算子实现。}
  \label{fig:pipeline}
\end{figure}

% =========================================================
\section{实验及分析}
\label{sec:exp}

\subsection{实验方法、工具、环境}

\paragraph{编程语言}
Python 3.12
\paragraph{主要依赖}
NumPy、SciPy、pandas、matplotlib、seaborn、joblib、scikit-posthocs、requests。

\paragraph{算法参数}
所有参数取文献推荐或常用默认值，统一种群规模约 50（仅 TSP GA 用 100，反映排列编码下
需要更大种群维持多样性），保证横向比较公平。具体参数见表~\ref{tab:params}。

\begin{table}[H]
  \centering
  \caption{六种算法形态的关键参数}
  \label{tab:params}
  \small
  \begin{tabularx}{\linewidth}{lX}
    \toprule
    Variant & Key parameters \\
    \midrule
    GA (continuous)   & pop $=50$, $p_c=0.9$, $\eta_c=15$, $\eta_m=20$ \\
    PSO (continuous)  & swarm $=50$, $w\!:\!0.9\!\to\!0.4$, $c_1=c_2=2.0$, $v_{\max}=0.2\,D$ \\
    ACOR (continuous) & $k=50$, $m=10$, $q=10^{-4}$, $\xi=0.85$ \\
    GA (TSP)          & pop $=100$, $p_c=0.9$, $p_m=0.2$, $t=3$ \\
    ACO (TSP)         & $m=n$, $\alpha=1$, $\beta=2$, $\rho=0.5$, $Q=1$ \\
    PSO (TSP)         & swarm $=50$, $w=0.3$, $c_1=0.7$, $c_2=0.9$ \\
    \bottomrule
  \end{tabularx}
\end{table}

\paragraph{符号说明}
\begin{itemize}\setlength{\itemsep}{0pt}
  \item pop / swarm —— 种群规模
  \item $p_c$ —— 交叉概率
  \item $p_m$ —— 变异概率
  \item $\eta_c$ —— SBX 交叉分布指数
  \item $\eta_m$ —— 多项式变异分布指数
  \item $w$ —— 惯性权重（连续 PSO 由 0.9 线性衰减至 0.4）
  \item $c_1,\ c_2$ —— 认知 / 社会加速系数
  \item $v_{\max}$ —— 速度上限；连续 PSO 取 $0.2\,D$，$D$ 为定义域宽度
  \item $k$ —— ACOR 解档案大小
  \item $m$ —— 蚂蚁数（TSP ACO 中 $m=n$，$n$ 为城市数）
  \item $q,\ \xi$ —— ACOR 局部性 / 收敛速率控制
  \item $\alpha,\ \beta$ —— 信息素与启发式指数
  \item $\rho$ —— 信息素蒸发率
  \item $Q$ —— 信息素沉积常数
  \item $t$ —— 锦标赛规模
  \item 算子：连续 GA 用 SBX + 多项式变异；TSP GA 用 OX 交叉 + 反转变异；TSP PSO 用 Clerc 2004 交换序列离散形式
\end{itemize}

\paragraph{实验规模}
\begin{itemize}\setlength{\itemsep}{0pt}
  \item 仿真 —— 3 算法 $\times$ 5 函数 $\times$ 15 重复 $= 225$ 次独立运行，维度 $d=30$，每次 FE 预算 30000
  \item 真实 TSP —— 3 算法 $\times$ 4 实例 $\times$ 15 重复 $= 180$ 次运行，每次 FE 预算 20000
\end{itemize}


\paragraph{评价与统计检验}
仿真主指标：最终最优 $f_{\rm best}$、成功率（$f<10^{-2}$）、收敛曲线；TSP 主指标：
最优巡回长度、gap\%、收敛曲线与最优路线图。实例内三算法两两采用 Wilcoxon 符号秩配对
检验（按种子配对，$\alpha=0.05$）；跨实例综合排名采用 Friedman 秩检验，
若整体显著则继续 Nemenyi 事后两两比较。

\subsection{仿真实验结果与分析}

\paragraph{主结果表}
表~\ref{tab:func_main} 汇总三算法在 5 个测试函数上的 15 次重复 best 值（均值 $\pm$ 标准差），
最优值（即该行三算法的最小均值）在数值后以上标 $\ast$ 标记，表注注明。

\begin{table}[H]
  \centering
  \caption{Appendix C 测试函数最终 best $f$ 均值 $\pm$ 标准差（$d=30$，15 次重复，越小越好；标 $\ast$ 者为该行各算法的最小值）}
  \label{tab:func_main}
  \small
  \begin{tabular}{lccc}
    \toprule
    函数 & GA & ACO（ACOR） & PSO \\
    \midrule
    Sphere     & $3.0\times 10^{-3} \pm 1.6\times 10^{-3}$
               & $7.8\times 10^{-40} \pm 1.7\times 10^{-39}\,^{\ast}$
               & $6.0\times 10^{-5} \pm 4.1\times 10^{-5}$ \\
    Rosenbrock & $44.2 \pm 26$
               & $20.7 \pm 15\,^{\ast}$
               & $27.0 \pm 0.9$ \\
    Rastrigin  & $21.1 \pm 4.3\,^{\ast}$
               & $125.3 \pm 29$
               & $57.0 \pm 20$ \\
    Ackley     & $0.83 \pm 0.37$
               & $9.38 \pm 5.0$
               & $0.33 \pm 0.56\,^{\ast}$ \\
    Griewank   & $0.80 \pm 0.13$
               & $0.17 \pm 0.48$
               & $0.07 \pm 0.06\,^{\ast}$ \\
    \bottomrule
  \end{tabular}
\end{table}

\paragraph{成功率}
将 $f_{\rm best}<10^{-2}$ 视为"成功"。除单峰 Sphere 三算法 100\% 成功外，
30 维下其余多峰函数普遍难以达标，仅 ACOR 在 Griewank 上有 27\% 成功率
（见表~\ref{tab:func_succ}）。

\begin{table}[H]
  \centering
  \caption{15 次重复中达成阈值 $f<10^{-2}$ 的比例}
  \label{tab:func_succ}
  \small
  \begin{tabular}{lccc}
    \toprule
    函数 & GA & ACO（ACOR） & PSO \\
    \midrule
    Sphere     & 1.00 & 1.00 & 1.00 \\
    Rosenbrock & 0.00 & 0.00 & 0.00 \\
    Rastrigin  & 0.00 & 0.00 & 0.00 \\
    Ackley     & 0.00 & 0.00 & 0.00 \\
    Griewank   & 0.00 & 0.27 & 0.00 \\
    \bottomrule
  \end{tabular}
\end{table}

\begin{figure}[H]
  \centering
  \includegraphics[width=0.95\linewidth]{fig1_func_convergence.pdf}
  \caption{5 个测试函数上三算法 best-so-far 平均收敛曲线（横轴函数评估次数，纵轴对数刻度）。}
  \label{fig:func_conv}
\end{figure}

\begin{figure}[H]
  \centering
  \includegraphics[width=0.92\linewidth]{fig2_func_box.pdf}
  \caption{5 函数最终最优值的箱线图（15 次重复，纵轴对数刻度）。}
  \label{fig:func_box}
\end{figure}

\paragraph{显著性}
跨 5 函数的 Friedman 秩检验得 $\chi^2=1.6$、$p=0.45$（未达 $0.05$ 显著水平），
平均秩 PSO $1.6 <$ ACO $2.0 <$ GA $2.4$。整体未显著反映出三算法在不同地形上互有
胜负、不存在普适赢家。两两 Wilcoxon 在大多数函数上仍能得到 $p<0.05$ 的实例内
显著差异。

\paragraph{结果分析}

\begin{enumerate}
  \item ACOR 在光滑地形上精度极高：Sphere 上达 $\sim10^{-40}$，比 PSO 的 $10^{-5}$
        与 GA 的 $10^{-3}$ 高出数十个数量级；谷型 Rosenbrock 上也最优（20.7）。
        这印证了 ACOR 的解档案 $+$ 高斯核采样在平滑、强单峰地形上的快速收敛优势。
  \item ACOR 在强多峰地形上明显劣化：Rastrigin 高达 125.3、Ackley 为 9.4，反映出
        档案过快收紧导致样本集中于局部极小区域、丧失全局探索能力，体现典型的"早熟"。
  \item PSO 在多峰地形上最强：Ackley（0.33）与 Griewank（0.07）均居首。
        群体多样性与惯性权重线性衰减的探索--开发平衡机制契合多峰场景。
  \item GA 在最崎岖的 Rastrigin 上反而最优（21.1）：SBX 与多项式变异结合提供了较强
        局部精修与广域扰动的平衡，使其在规则多峰陷阱密集的地形中维持较好的种群多样性。
  \item Friedman 平均秩 PSO $<$ ACO $<$ GA 表明 PSO 在连续测试函数集合上整体最好，
        但三者秩差较小，体现"地形决定胜负"的 No Free Lunch 现象。
\end{enumerate}

\subsection{真实数据（TSP）实验结果与分析}

\paragraph{主结果表}
表~\ref{tab:tsp_main} 汇总三算法在 4 个 TSP 实例上的 gap\%（越小越好），
标 $\ast$ 者为该行最小值。

\begin{table}[H]
  \centering
  \caption{TSP 各实例 15 次重复 gap\% 均值 $\pm$ 标准差（越小越好；标 $\ast$ 者为该行各算法的最小值）}
  \label{tab:tsp_main}
  \small
  \begin{tabular}{lccc}
    \toprule
    实例 & GA & ACO & PSO \\
    \midrule
    berlin52 & $46.6 \pm 6.4$  & $0.9 \pm 0.8\,^{\ast}$ & $175.4 \pm 17.0$ \\
    eil51    & $53.6 \pm 6.3$  & $5.2 \pm 1.4\,^{\ast}$ & $163.7 \pm 10.3$ \\
    kroA100  & $203.0 \pm 9.7$ & $8.1 \pm 1.2\,^{\ast}$ & $452.7 \pm 23.1$ \\
    ch150    & $267.6 \pm 8.8$ & $6.5 \pm 1.4\,^{\ast}$ & $507.2 \pm 25.5$ \\
    \bottomrule
  \end{tabular}
\end{table}

\begin{figure}[H]
  \centering
  \includegraphics[width=0.92\linewidth]{fig3_tsp_box.pdf}
  \caption{4 个 TSP 实例的 gap\% 箱线图（15 次重复）。}
  \label{fig:tsp_box}
\end{figure}

\begin{figure}[H]
  \centering
  \includegraphics[width=0.88\linewidth]{fig4_tsp_convergence.pdf}
  \caption{kroA100 上三算法的 best-so-far 巡回长度收敛曲线（均值 $\pm$ 标准差，阴影为 $\pm 1$ 标准差带）。}
  \label{fig:tsp_conv}
\end{figure}

\begin{figure}[H]
  \centering
  \includegraphics[width=0.95\linewidth]{fig5_tsp_routes.pdf}
  \caption{berlin52 上三算法的最优巡回路线对比（每子图取该算法 15 个种子中找到的最短巡回）。}
  \label{fig:tsp_routes}
\end{figure}

\paragraph{显著性}
跨 4 实例 Friedman 秩检验 $\chi^2=8.0$、$p=0.018$（在 $\alpha=0.05$ 下显著），
平均秩 ACO $1.0 <$ GA $2.0 <$ PSO $3.0$。Nemenyi 事后检验中 ACO 显著优于 PSO
（$p=0.013$）；ACO 与 GA、GA 与 PSO 在 $\alpha=0.05$ 下未达 Nemenyi 显著，但配对
Wilcoxon 全部 $p<10^{-3}$。

\paragraph{结果分析}

\begin{enumerate}
  \item ACO 在 TSP 上绝对占优：berlin52 上 gap 仅 0.9\%，几乎达到已知最优 7542；
        其余实例 gap 均在 10\% 以内。信息素 $+$ 启发式 $\eta=1/d_{ij}$ 的协同构造直接
        利用了 TSP "短边优先连接"的结构，是该算法的本职舞台。图~\ref{fig:tsp_routes}
        中 ACO 的巡回光滑且无明显交叉，与另两算法形成强烈对比。
  \item 离散 PSO 表现最差：gap 普遍在 160\%--500\% 之间。连续 PSO 的"速度向量"
        概念在排列空间中方向性弱化为交换序列，社会与认知系数对位置的引导大幅退化，
        巡回杂乱多交叉。这与文献对离散 PSO 在 TSP 上的报告完全一致。
  \item GA 居中：OX 与反转变异具备基本的局部改良能力，但缺少 TSP 专用的局部搜索
        （如 2-opt），规模扩大后 gap 显著上升（kroA100 达 203\%，ch150 达 268\%）。
  \item ACO 的方差最小（多数 $<2$，远小于另两算法的 $10$ 量级），说明其多次重复
        运行的结果稳健、可重复。
\end{enumerate}

\subsection{综合比较}

将仿真与真实两类场景的 Friedman 平均秩并列汇总于表~\ref{tab:overall}。

\begin{table}[H]
  \centering
  \caption{两场景综合排名（按 Friedman 平均秩，1 = 最优；标 $\ast$ 者为该场景秩最小）}
  \label{tab:overall}
  \small
  \begin{tabular}{lcccl}
    \toprule
    场景 & GA & ACO & PSO & 显著性 \\
    \midrule
    连续测试函数（5 实例） & 2.4 & 2.0 & $1.6\,^{\ast}$ & $\chi^2=1.6,\ p=0.45$（不显著） \\
    真实 TSP（4 实例）     & 2.0 & $1.0\,^{\ast}$ & 3.0 & $\chi^2=8.0,\ p=0.018$（显著） \\
    \bottomrule
  \end{tabular}
\end{table}

\paragraph{核心规律}
ACO 在本职的组合问题 TSP 上以绝对优势占第一名（Friedman 显著，Nemenyi
ACO $\succ$ PSO $p=0.013$）；但其连续推广 ACOR 在多峰函数上反而最差，
说明算法的优势高度依赖于问题结构契合度。PSO 在连续场景平均秩最优，但在组合场景沦为
最差，体现"速度概念"在排列空间难以保留。GA 在两类场景下均居中，是稳健的全能选手。
这正是 No Free Lunch 定理的直接实证：不存在对所有问题都最优的算法，
算法选型必须匹配问题的编码类型与地形特性。

% =========================================================
\section{结论}
\label{sec:conclusion}

\begin{enumerate}
  \item 三种仿生算法在两类问题上均能正常工作并明显优于随机搜索：连续 Sphere 上三算法
        均成功收敛到 $f<10^{-2}$；真实 TSP 上各算法均能产出有效巡回，证明仿生优化在实际
        场景中的可行性与基本有效性。
  \item ACO 在 TSP 上以绝对优势夺冠（berlin52 gap 仅 0.9\%，Friedman 平均秩 1.0，
        Nemenyi 显著优于 PSO），体现了"算法--问题"结构契合的重要性 ——
        ACO 设计之初就是为 TSP 而生，其信息素引导与启发式构造直接对应 TSP 的最短边偏好结构。
  \item ACO 的连续推广 ACOR 在光滑单峰（Sphere、Rosenbrock）上精度极高，但在强多峰
        （Rastrigin、Ackley）上明显早熟最差，说明跨问题域外推时算法的强项与短板都会被放大。
  \item PSO 在连续多峰问题上整体最好（Friedman 平均秩 1.6），但离散化到 TSP 后退化为最弱
        （平均秩 3.0），暴露了基于速度向量的算法在组合空间中的根本困难。
  \item GA 在两个场景下均居中，没有明显的天花板，但也没有突出的强项 ——
        作为重组型通用框架，它的"全能"也意味着"在任何具体问题上都不是第一名"。
  \item No Free Lunch 定理在本报告中得到完整实证：ACO 与 PSO 在不同场景互换冠/末位、
        GA 始终居中。在真实选型时应先分析问题结构（连续或组合、单峰或多峰、维度高低、
        是否带约束），再匹配相应算法或其专用变体。
\end{enumerate}

% =========================================================
\appendix
\section*{附录\ \ 关键公式}
\addcontentsline{toc}{section}{附录\ \ 关键公式}

\paragraph{Appendix C 测试函数（维度 $d$，最小化，全局最优 $f^{\ast}=0$）。}
\begin{align*}
  f_{\rm Sphere}(\bm{x})     &= \sum_{i=1}^{d} x_i^2, \\
  f_{\rm Rosenbrock}(\bm{x}) &= \sum_{i=1}^{d-1}\!\Big[\,100\,(x_{i+1}-x_i^{\,2})^2+(x_i-1)^2\Big], \\
  f_{\rm Rastrigin}(\bm{x})  &= 10\,d + \sum_{i=1}^{d}\!\Big[\,x_i^{\,2}-10\cos(2\pi x_i)\Big], \\
  f_{\rm Ackley}(\bm{x})     &= -20\exp\!\bigg(\!-0.2\sqrt{\tfrac{1}{d}\textstyle\sum x_i^{\,2}}\bigg)
                               -\exp\!\bigg(\tfrac{1}{d}\textstyle\sum\cos(2\pi x_i)\bigg) + 20 + e, \\
  f_{\rm Griewank}(\bm{x})   &= \sum_{i=1}^{d}\frac{x_i^{\,2}}{4000} - \prod_{i=1}^{d}\cos\!\left(\frac{x_i}{\sqrt{i}}\right) + 1.
\end{align*}

\paragraph{TSP 巡回长度与 EUC\_2D 距离}
对排列 $\pi=(\pi_1,\dots,\pi_n)$，巡回长度为
\[
  L(\pi)=\sum_{i=1}^{n-1} d(\pi_i,\pi_{i+1}) + d(\pi_n,\pi_1).
\]
TSPLIB 约定的 EUC\_2D 距离：$d(u,v)=\mathrm{round}\!\big(\sqrt{(x_u-x_v)^2+(y_u-y_v)^2}\big)$。
相对已知最优的差距：$\mathrm{gap}\,\% = (L_{\rm best}-L^{\ast})/L^{\ast}\times 100\%$。

\paragraph{ACOR 高斯核采样（连续 ACO）}
按解档案中第 $l$ 个解的 rank 权重 $w_l$ 选定引导解，对每一维 $i$ 抽样
\[
  x_i \sim \mathcal{N}\!\Big(\mu_{l,i},\,\sigma_{l,i}^{\,2}\Big),\qquad
  \mu_{l,i} = s_{l,i},\qquad
  \sigma_{l,i} = \frac{\xi}{k-1}\sum_{e=1}^{k}\bigl|\,s_{e,i}-s_{l,i}\,\bigr|.
\]
新蚁与档案合并后按目标值排序、保留最优 $k$ 个，完成一次"信息素更新"。

\paragraph{Ant System 信息素更新（TSP ACO）}
所有蚂蚁完成巡回构造后，先按比率 $\rho$ 蒸发，再按巡回质量沉积：
\[
  \tau_{ij} \leftarrow (1-\rho)\,\tau_{ij} + \sum_{a=1}^{m}\Delta\tau_{ij}^{(a)},\qquad
  \Delta\tau_{ij}^{(a)} =
    \begin{cases}
      Q/L_a, & (i,j)\in T_a, \\
      0, & \text{其他.}
    \end{cases}
\]
蚂蚁 $a$ 从当前城市 $i$ 选择下一城市 $j$ 的概率：
\[
  p_{ij}^{(a)} \;\propto\; [\tau_{ij}]^{\alpha}\,[\eta_{ij}]^{\beta},\qquad \eta_{ij}=1/d_{ij}.
\]

\paragraph{交换序列离散 PSO（TSP PSO）}
将位置定义为排列 $X$，速度定义为有序交换算子序列
$V=[(i_1,j_1),(i_2,j_2),\dots]$。每代以惯性概率 $w$、认知概率 $c_1$、社会概率 $c_2$
分别把"$X$ 变成 $p_{\rm best}$ / $g_{\rm best}$ 所需的交换序列"按位概率施加到 $X$ 上：
\[
  X^{(t+1)} = X^{(t)} \oplus w\!\cdot\! V^{(t)}
              \oplus c_1\!\cdot\!(p_{\rm best}\!\ominus\!X^{(t)})
              \oplus c_2\!\cdot\!(g_{\rm best}\!\ominus\!X^{(t)}),
\]
其中 $\ominus$ 表示"从一个排列到另一个排列所需的交换序列"，$\oplus$ 表示"按概率
连续应用交换算子"。



\end{document}
